\documentclass[phy1200]{handout}

\title{Polarization measurements and \\superposition}

\begin{document}

\maketitle

Light has a polarization, meaning its electric field is pointing in a certain direction. Usually all the light rays emitted from an object have random polarizations.

However, there are special filters that selectively filter most polarizations. This means that the light going out of a polarization filter is \textbf{polarized}, meaning it only has a single polarization. 

But we can pass this polarized light through further filters. What happens to light as you pass light through multiple filters requires some careful explanation.

An ideal polarizing filter lets through 100\% of light with the polarization aligned to the filter and the light on the other side will only have a single polarization. In practice, though, filters will always absorb some of the light (even the part that is supposed to pass through completely). Additionally, some stray non-aligned polarization light will leak through the filter. In your explanations, you can pretend the filters are ideal.

\section{Two filters}
First, we will see what happens when we pass the light through two filters. Given a source beam of unpolarized light, the first filter simply polarizes the beam. 

Align the polarization axes of the two filters and observe the intensity of the beam. We will pretend that this is 100\% of the polarized beam.

Now rotate the second polarizer. Observe what happens and try to develop an explanation of what is going on (it should involve some trigonometric function). Try to explain why this would occur in terms of the polarization vector of the light.

\section{Three filters}
Set the two original filters to have their axes perpendicular to each other. Now place a third filter in between these two with its axis aligned to the first. Rotate it and observe the effect. 

Revise your explanation of the two-filter effect until you think that you can explain both cases in the same way.

Read or find videos about quantum superposition and spin measurements\footnote{One good resource is ``Energy and Matter: Our Quantum World (Knop)'' chapter four. This an open source textbook available for free online. }

\section{Report}
Things to include in your report.
\begin{enumerate}
    \item What is the polarization of a light beam?
    \item What does a polarizing filter do?
    \item Explain the behavior of the two-filter setup (with any necessary pictures of the beam on the screen).
    \item Explain the behavior of the three-filter setup (with any necessary pictures of the beam on the screen).
    \item Explain how this is related to the idea of `quantum superposition' and repeated spin measurements.    
\end{enumerate}




\end{document}